\section{R ile çözümlü uygulama}
\label{sec:r-bootstrap}

\textbf{Soru.} Dış veri dosyası gerektirmeyen iki öğretim örneği kullanın:
$2,3,3,4,13$ değerlerinin ortalaması için bootstrap standart hatası ve
aralığı; $7,8,9$ ile $2,3,4$ grupları için tam permütasyon testi hesaplayın.
Bunlar bölümdeki gerçek veri analizlerinin yerine geçen yeni bulgular değildir.

\begin{lstlisting}[style=prismR]
set.seed(2026)
degerler <- c(2, 3, 3, 4, 13)
tekrar <- 10000
bootstrap <- replicate(tekrar,
  mean(sample(degerler, length(degerler), replace = TRUE)))
c(ortalama = mean(degerler), bootstrap_se = sd(bootstrap))
quantile(bootstrap, c(0.025, 0.975), type = 7)
sqrt(mean((degerler - mean(degerler))^2) / length(degerler))
tam_bootstrap <- expand.grid(rep(list(degerler), length(degerler)))
quantile(rowMeans(tam_bootstrap), c(0.025, 0.975), type = 7)
grup_a <- c(7, 8, 9)
grup_b <- c(2, 3, 4)
birlesik <- c(grup_a, grup_b)
gozlenen <- mean(grup_a) - mean(grup_b)
atamalar <- combn(seq_along(birlesik), length(grup_a))
farklar <- apply(atamalar, 2, function(indis) {
  mean(birlesik[indis]) - mean(birlesik[-indis])
})
p_tam <- mean(abs(farklar) >= abs(gozlenen) - 1e-12)
c(fark = gozlenen, atama_sayisi = ncol(atamalar), p = p_tam)
\end{lstlisting}

\begin{outputguidebox}
\textbf{Çözüm ve kontrol değerleri.} İlk örneğin ortalaması 5'tir.
Benzetimle bulunan bootstrap standart hatası, ampirik dağılım altında
hesaplanan 1,8111'e yaklaşır. Bütün $5^5=3125$ geri koymalı sıralı
örneklem listelenirse tam bootstrap yüzdelik aralığı $[2{,}6;9{,}0]$ olur;
10.000 tekrarlı benzetimin uçları küçük farklılıklar gösterebilir.
İkinci örnekte gözlenen fark $8-3=5$'tir. Olası $\binom{6}{3}=20$
atamanın ikisi mutlak değerce en az bu kadar uçtur; tam iki yönlü $p=0{,}10$ olur.
\end{outputguidebox}

\textbf{Yorum.} Bootstrap geri koymalı, grup etiketlerinin permütasyonu
geri koymasızdır. Yalnız beş gözlem ve belirgin bir uç değer içeren ilk örnek,
güvenilir kapsama garantisi değil yöntem gösterimidir. Tam sayımda
\texttt{+1} düzeltmesi uygulanmaz; bu düzeltme sınırlı rastgele permütasyonla
yaklaştırılan testler içindir. Permütasyon testi uygun etiket değiştirilebilirliği
veya bu atamaları üreten rastgele atama tasarımını gerektirir.

\textbf{Pekiştirme ve yanıt.} Gözlenen fark büyük görünse de bu çok küçük
örnekte çift yönlü $p=0{,}10$ olduğundan yüzde 5 düzeyinde $H_0$ reddedilemez.
Bu sonuç iki evrenin aynı olduğunun kanıtı değildir.